
\section{Funtores exactos entre categorías galoisianas}
\label{V.6}

\begin{proposition}
\label{V.6.1}
Sean $\cal{C}$ y $\cal{C'}$ dos categorías galoisianas,
$H\colon \cal{C}\to \cal{C'}$ un funtor covariante, $F'$ un funtor
fundamental sobre~$\cal{C'}$ y $F=F'\circ H$. Son equivalentes las
condiciones siguientes:
\begin{enumerate}
\item[(i)] $H$ es \emph{exacto}, \ie exacto por la izquierda y exacto
por la derecha.
\item[(ii)] $H$ es exacto por la izquierda, transforma sumas finitas en
sumas finitas y epimorfismos en epimorfismos (o, equivalentemente,
objetos $\not\approx\emptyset_{\cal C}$ en objetos
$\not\approx\emptyset_{\cal C'}$).
\item[(iii)] $F$ es un funtor fundamental sobre~$\cal C$.
\end{enumerate}
\end{proposition}

La implicación (i)$\To$(ii) es un hecho general de la teoría de categorías. Además,
la primera forma dada a (ii) implica la segunda, como se ve observando que,
si~$X$ es un objeto de~$\cal C$, entonces
$X\not\approx\emptyset_{\cal C}$ \sss\ el morfismo
$X\to e_{\cal C}$ es un epimorfismo; se observará que~$F$, al suponerse
exacto por la izquierda, transforma $e_{\cal C}$ en~$e_{\cal C'}$. La
segunda forma dada a (ii) implica (iii), pues~$F$, al ser exacto por la
izquierda y, por tanto, prorepresentable, satisface el criterio (iii)
de~\Ref{V.5.6}. Por último, (iii) implica (i), como se sigue de que~$F$ es
exacto y ``conservativo'' (\ie satisface el axioma (G~6) del
\No~\Ref{V.4}).

Sea entonces~$\Gamma$ el grupoide fundamental de~$\cal C$ y~$\Gamma'$ el
de~$\cal C'$. Así, si
\marginpar{135}
$H$ es exacto, entonces $F'\mto F'\circ H$ es un funtor del
grupoide~$\Gamma'$ en el grupoide~$\Gamma$, que denotaremos por~$\tH$:
$$
\tH(F')(X)=F'(H(X)).
$$
También puede escribirse, con la notación $F(X)=E_X(F)$ introducida en el
\ifthenelse{\boolean{orig}}
{\No~\Ref{V.6}:}
{\No~\Ref{V.5}:}
$$
\ifthenelse{\boolean{orig}}
{E_{H(X)}(F')=E_X(\tH(F'))}
{E_{H(X)}(F')=E_X(\tH(F')).}
$$
Esta última fórmula muestra, teniendo en cuenta~\Ref{V.5.8}
o~\Ref{V.4.1}, que el funtor exacto~$H$ queda determinado (salvo
isomorfismo único) cuando se conoce el funtor correspondiente~$\tH$.
Fijemos un~$F'$ y sea $F=\tH(F')$; entonces~$\tH$ define un homomorfismo
de~$\pi_{F'}=\Aut(F')$ en~$\pi_F=\Aut(F)$:
$$
\tH\colon\pi_{F'}\to\pi_F
\qquad(F=\tH(F')=F'\circ H).
$$
Por otra parte, la fórmula anterior muestra (teniendo en cuenta~\Ref{V.5.8})
que este homomorfismo tiene la propiedad siguiente: para todo conjunto
finito~$E$ sobre el que~$\pi_F$ actúa continuamente, el grupo~$\pi_{F'}$
actúa también \emph{continuamente} por medio del homomorfismo anterior
$\pi_{F'}\to\pi_F$. Aplicándolo a los cocientes de~$\pi_F$ por sus
subgrupos normales abiertos, se ve que la condición anterior significa
también que el homomorfismo considerado es continuo. Recíprocamente,
démonos un objeto~$F$ de~$\Gamma$, un objeto~$F'$ de~$\Gamma'$ y un
homomorfismo continuo
$$
u\colon\pi_{F'}\to\pi_F.
$$
Le corresponde entonces un funtor de~$\cal C(\pi)$ en~$\cal C(\pi')$,
manifiestamente exacto; por tanto, en virtud de~\Ref{V.4.1}, le corresponde
un funtor $H$ de~$\cal C$ en~$\cal C'$ que es exacto y tal que
$\tH\colon\pi_{F'}\to\pi_F$ es precisamente~$u$. En lugar de partir de un
homomorfismo de grupos, también puede partirse de un \emph{funtor}
$$
U\colon\Gamma'\to\Gamma
$$
tal que, para \emph{todo} $F'\in\Gamma'$ (o para \emph{un}
$F'\in\Gamma'$, lo que es lo mismo), el homomorfismo correspondiente
$\pi_{F'}\to\pi_F$ sea continuo: tal funtor es isomorfo a un funtor de la
forma~$\tH$, donde $H\colon\cal C\to\cal C'$ es un funtor exacto
determinado salvo isomorfismo único. Así:

\begin{corollaire}
\label{V.6.2}
Para
\marginpar{136}
que un funtor $H\colon\cal C\to\cal C'$ de categorías galoisianas sea
exacto, es necesario y suficiente que existan equivalencias
$\cal C(\pi)\to\cal C$ y
\ifthenelse{\boolean{orig}}
{$\cal C'\to\cal C(\pi)$}
{$\cal C'\to\cal C(\pi')$}
que transformen el funtor~$H$ en el funtor
$\cal C(\pi)\to\cal C(\pi')$ asociado a un homomorfismo de grupos
topológicos $\pi'\to\pi$.
\end{corollaire}

\begin{corollaire}
\label{V.6.3}
Sean $\cal C$ y $\cal C'$ dos categorías galoisianas y sean~$\Gamma$ y
$\Gamma'$ sus grupoides fundamentales. Entonces la categoría de los
funtores exactos de~$\cal C$ en~$\cal C'$ es equivalente a la categoría de
los funtores $U\colon\Gamma'\to\Gamma$ que tienen la propiedad siguiente:
para todo $F'$ de~$\Gamma'$ (o para \emph{un} $F'$ de~$\Gamma'$, lo que es
lo mismo), poniendo $F=U(F')$, el homomorfismo
$$
\pi_{F'}=\Aut(F')\to\pi_F=\Aut(F)
$$
definido por~$U$ es continuo.
\end{corollaire}

Consideremos el progrupo fundamental~$\Pi$ de~$\cal C$. Entonces un funtor
exacto~$H$ lo transforma en un progrupo $H(\Pi)$ de~$\cal C'$. Vamos a
definir un homomorfismo
$$
\Pi'\to H(\Pi)
$$
(donde~$\Pi'$ es el progrupo fundamental de~$\cal C'$), mediante la
condición de que, para todo objeto~$F'$ de~$\Gamma'$, el homomorfismo
correspondiente
$$
F'(\Pi')=\pi_{F'}\to F'(H(\Pi))=\pi_F\qquad
\ifthenelse{\boolean{orig}}
{\text{(donde $F=F'\circ H=\tH(F')$)};}
{\text{(donde $F=F'\circ H=\tH(F')$)}}
$$
sea el homomorfismo natural
$$
\Aut(F')\to\Aut(F'\circ H).
$$
Como este último es funtorial en~$F'$, define efectivamente, en virtud
de~\Ref{V.5.8}, un homomorfismo de proobjetos, y de hecho de progrupos,
en~$\cal C'$. Este homomorfismo se llama \emph{asociado} al funtor~$H$.

Sea ahora~$H'$ un segundo funtor exacto, de la categoría galoisiana
$\cal C'$ en una categoría galoisiana~$\cal C''$. Es inmediato que
$$
{}^t(H'H)=\tH\,\tH'
$$
(N.B.: se trata de una identidad de funtores, y no solamente de un
isomorfismo canónico). Hay una propiedad de transitividad análoga para los
homomorfismos asociados de los progrupos fundamentales.

Vamos
\marginpar{137}
a interpretar ahora las propiedades del funtor exacto~$H$ en términos del
homomorfismo correspondiente
$$
u\colon\pi_{F'}\to\pi_F\qquad
\ifthenelse{\boolean{orig}}
{\text{(donde $F'=F'\circ H$)}.}
{\text{(donde $F=F'\circ H$)}.}
$$
Conviene introducir la noción de \emph{objeto punteado}
\index{punteado (objeto)|hyperpage}%
de la categoría galoisiana~$\cal C$ (provista de su funtor fundamental~$F$):
es, por definición, un objeto~$X$ de~$\cal C$ provisto de un elemento~$a$
de~$F(X)$. Se interpreta, pues, como un conjunto finito sobre el que
$\pi_F$ actúa continuamente por la izquierda, provisto de un punto~$a$.
Por tanto, los objetos punteados \emph{conexos} de~$\cal C$ se identifican,
en virtud de~\Ref{V.5.3}, con los subgrupos abiertos de~$\pi_F$. Si~$U$ y
$V$ son dos de tales subgrupos, correspondientes a objetos punteados
conexos~$X$ e~$Y$ de~$\cal C$, entonces existe un homomorfismo punteado
de~$X$ en~$Y$ si y solo si $U\subset V$, y dicho homomorfismo es entonces
único. Por supuesto, el funtor~$H$ transforma objetos punteados en objetos
punteados (pues $F=F'\circ H$). Obsérvese, por otra parte, que un subgrupo
cerrado de un grupo como~$\pi_F$ es la intersección de los subgrupos
abiertos que lo contienen; por consiguiente,
\ifthenelse{\boolean{orig}}{\ignorespaces}{si}
$M$ y~$N$ son dos subgrupos cerrados, entonces $M\subset N$ si y solo si
todo subgrupo abierto que contiene a~$N$ contiene también a~$M$. Gracias a
estas observaciones se demuestran fácilmente los resultados siguientes:
